\exercisesection{2}

\begin{enumerate}
\def\labelenumi{\arabic{enumi}.}
\item
  Let K be a commutative field, A the polynomial ring K{[}X, Y{]} in two indeterminates and \(a_{n}\) the principal ideal \((XY^{n})\) in A; the sequence \((a_{n})_{n\geqslant1}\) forms with \(a_{0}=A\) an exhaustive and separated filtration on A. Let b be the principal ideal (X) of A; show that the topology on b induced by that on A is strictly coarser than the topology defined by the filtration on b associated with that on A (and a fortiori this latter filtration is distinct from the filtration induced by that on A).
\item
  Let K be a commutative field of characteristic \(\neq2\) and A the ring K{[}{[}X, Y{]}{]} of formal power series in two indeterminates.
\end{enumerate}

\begin{enumerate}
\def\labelenumi{(\alph{enumi})}
\item
  Show that in A the principal ideal \(\mathfrak{p} = (\mathrm{X}^2 -\mathrm{Y}^3)\) is prime. (If a product \(f(\mathbf{X},\mathbf{Y})g(\mathbf{X},\mathbf{Y})\) of two formal power series is divisible by \(\mathbf{X}^2 -\mathbf{Y}^3\) , note first that \(f(\mathrm{T}^3,\mathrm{T}^2) = 0\) or \(g(\mathrm{T}^3,\mathrm{T}^2) = 0\) in the ring of formal power series \(\mathbf{K}[[\mathrm{T}]]\) ; assuming for example \(f(\mathrm{T}^3,\mathrm{T}^2) = 0\) , show first that \(f(\mathbf{X},\mathbf{Y}^2)\) is divisible by \(\mathbf{X} - \mathbf{Y}^3\) and by \(\mathbf{X} + \mathbf{Y}^3\) and, by considering \(f(-\mathbf{Y}^3 +\mathbf{X},\mathbf{Y}^2)\) , prove finally that \(f(\mathbf{X},\mathbf{Y}^2)\) is divisible by \(\mathbf{X}^2 -\mathbf{Y}^6\) .)
\item
  Let \(m\) be the maximal ideal \(AX + AY\) of \(A\) . Show that \(p\) is closed with respect to the \(m\) -adic topology on \(A\) (with which \(A\) is Hausdorff and complete) but that in the ring \(\text{gr}(A)\) , \(\text{gr}(p)\) is not a prime ideal ( \(p\) being given the filtration induced by that on \(A\) ).
\end{enumerate}

\begin{enumerate}
\def\labelenumi{\arabic{enumi}.}
\setcounter{enumi}{2}
\item
  Let A be a filtered ring and E a finitely generated A-module. Show that, if E is given the filtration induced by that on A, gr(E) is a finitely generated gr(A)-module (cf.~Exercise 5 (c)).
\item
  Give an example of a bijective A-linear mapping \(u \colon \mathbf{E} \to \mathbf{F}\) , where \(\mathbf{E}\) and \(\mathbf{F}\) are two filtered A-modules, such that \(u\) is compatible with the filtrations but \(\operatorname{gr}(u)\) is neither injective nor surjective. (Take \(\mathbf{E} = \mathbf{A}\) , \(\mathbf{F} = \mathbf{A}\) with another filtration and \(u\) the identity mapping.)
\item
  Let A be a filtered commutative ring with filtration \((\mathfrak{a}_n)_{n\geqslant 0}\) such that \(\mathfrak{a}_0 = \mathbf{A}\) .
\end{enumerate}

\begin{enumerate}
\def\labelenumi{(\alph{enumi})}
\tightlist
\item
  For the ring gr(A) to be generated by a family of elements whose degrees are bounded, it is necessary and sufficient that there exist an integer q such that, for all \(n\) , \(\mathfrak{a}_n = \mathfrak{a}_{n+1} + \mathfrak{b}_n\) , where \(\mathfrak{b}_n\) is the sum of the ideals \(\mathfrak{a}_1^{\alpha_1}\mathfrak{a}_2^{\alpha_2}\ldots\mathfrak{a}_q^{\alpha_q}\) for
\end{enumerate}

\[
\alpha_ {i} \geqslant 0 (1 \leqslant i \leqslant q)
\]

\[
\sum_ {i = 1} ^ {q} i \alpha_ {i} = n.
\]

\[
\mathfrak {a} _ {n} = \mathfrak {a} _ {n + k} + \mathfrak {b} _ {n}
\]

\[
k > 0
\]

\begin{enumerate}
\def\labelenumi{(\alph{enumi})}
\setcounter{enumi}{1}
\item
  For \(\mathrm{gr}(\mathbf{A})\) to be Noetherian, it is necessary and sufficient that \(\mathbf{A} / \mathfrak{a}_1\) be Noetherian, that the condition of (a) hold and that for \(i \leqslant q\) the \(\mathfrak{a}_i / \mathfrak{a}_{i+1}\) be finitely generated \((\mathbf{A} / \mathfrak{a}_1)\) -modules (use Corollary 4 to Theorem 2 of no. 10).
\item
  Let K be a commutative field and A the polynomial ring in two indeterminates K{[}X, Y{]}. A total ordering is defined on the set of monomials \(X^{m}Y^{n}\) by setting \(X^{m}Y^{n} \leqslant X^{p}Y^{q}\) if \(m + n \leqslant p + q\) or if \(m + n = p + q\) and \(m \leqslant p\) . Let \((\mathbf{M}_{k})\) be the sequence of monomials in X, Y thus arranged in ascending order and let \(a_{n}\) be the ideal of A generated by the \(M_{k}\) of index \(k \geqslant n\) . Show that \((a_{n})\) is a filtration compatible with the ring structure on A; with this filtration gr(A) admits divisors of 0 and is not Noetherian, although A is a Noetherian integral domain (use the criterion in (b)); in particular, gr( \(a_{1}\) ) (with the filtration induced on \(a_{1}\) by that on A) is not a finitely generated gr(A)-module, although \(a_{1}\) is a finitely generated A-module (cf.~§ 1, no. 2, Corollary to Proposition 1).
\end{enumerate}

¶ 6. Let A be a commutative ring, E and F two A-modules and (E \(_{n}\) ) (resp. (F \(_{n}\) )) an exhaustive filtration on E (resp. F) consisting of sub-A-modules. On the tensor product G = E ⊗A F, consider the exhaustive filtration consisting of the G \(_{n}\) = ∑ \(_{i+j=n}\) Im(E \(_{i}\) ⊗A F \(_{j}\) ).

\begin{enumerate}
\def\labelenumi{(\alph{enumi})}
\tightlist
\item
  Show that the composite canonical homomorphisms
\end{enumerate}

\[
\begin{array}{r l}&(\mathrm{E} _ {i} / \mathrm{E} _ {i + 1}) \otimes_ {\mathrm{A}} (\mathrm{F} _ {j} / \mathrm{F} _ {j + 1}) \rightarrow\\&\qquad (\mathrm{E} _ {i} \otimes_ {\mathrm{A}} \mathrm{F} _ {j}) / (\operatorname{Im} (\mathrm{E} _ {i} \otimes_ {\mathrm{A}} \mathrm{F} _ {j + 1}) + \operatorname{Im} (\mathrm{E} _ {i + 1} \otimes_ {\mathrm{A}} \mathrm{F} _ {j})) \rightarrow \mathrm{G} _ {i + j} / \mathrm{G} _ {i + j + 1}\end{array}
\]

are the restrictions of a graded homomorphism of degree 0 (called canonical) \(\operatorname{gr}(\mathbf{E}) \otimes_{\mathbf{A}} \operatorname{gr}(\mathbf{F}) \to \operatorname{gr}(\mathbf{E} \otimes_{\mathbf{A}} \mathbf{F})\) , which is surjective.

\begin{enumerate}
\def\labelenumi{(\alph{enumi})}
\setcounter{enumi}{1}
\item
  Show that, if \(\operatorname{gr}(\mathbf{E})\) is a flat A-module, the A-modules \(\mathrm{E}_m / \mathrm{E}_n\) for \(m \leqslant n\) and \(\mathrm{E} / \mathrm{E}_n\) for \(n \in \mathbf{Z}\) are flat.
\item
  Suppose in the following that \(\mathrm{gr}(\mathbf{E})\) is a flat A-module and that \(\mathbf{E}\) is a flat A-module (the second hypothesis being a consequence of the first if \((\mathbf{E}_n)\) is a discrete filtration). Then the \(\mathrm{H}_{ij} = \mathrm{E}_i \otimes_{\mathbf{A}} \mathrm{F}_j\) are canonically identified with submodules of \(\mathbf{G} = \mathbf{E} \otimes_{\mathbf{A}} \mathbf{F}\) (Chapter I, § 2, no. 5, Proposition 4). Show that, for any two finite subsets \(\mathbf{R}\) , \(\mathbf{S}\) of \(\mathbf{Z} \times \mathbf{Z}\) ,
\end{enumerate}

\[
\left(\sum_ {(i, j) \in \mathbf {R}} \mathrm{H} _ {i j}\right) \cap \left(\sum_ {(h, k) \in \mathbf {S}} \mathrm{H} _ {h k}\right) = \mathrm{H} _ {\sup (i, h), \sup (j, k)}\tag{*}
\]

where in the second sum \((i,j,h,k)\) run through \(R \times S\) . (If R and S are each reduced to a single element, use Chapter I, §2, no. 6, Proposition 7. If p is the greatest of the indices i and h appearing in the elements of R or S, q the least of the indices j and k appearing in the elements of R or S, consider the images of the two sides of (*) under the canonical homomorphism

\[
\mathbf {E} \otimes_ {\mathbf {A}} \mathbf {F} _ {q} \rightarrow (\mathrm{E} / \mathrm{E} _ {p}) \otimes_ {\mathbf {A}} \mathbf {F} _ {q}
\]

and argue by induction on \(\operatorname{Card}(\mathbf{R}) + \operatorname{Card}(\mathbf{S})\) .

\begin{enumerate}
\def\labelenumi{(\alph{enumi})}
\setcounter{enumi}{3}
\tightlist
\item
  Deduce from (c) that the canonical homomorphism defined in (a) is then bijective.
\end{enumerate}

\begin{enumerate}
\def\labelenumi{\arabic{enumi}.}
\setcounter{enumi}{6}
\tightlist
\item
  Let A be a commutative ring, E an A-module, F a submodule of E, B the exterior algebra \(\wedge\) (E) and \(\Im\) the two-sided ideal of B generated by the canonical images in B of the elements of F. Let \(\operatorname{gr}^{\Im}(B)\) be the graded ring associated with the ring B filtered by the \(\Im\) -adic filtration.
\end{enumerate}

\begin{enumerate}
\def\labelenumi{(\alph{enumi})}
\item
   Define a surjective canonical (non-graded) A-algebra homomorphism \((\wedge (\mathbf{F})) \otimes_{\mathbf{A}}^{g}(\wedge (\mathbf{E} / \mathbf{F})) \to \mathrm{gr}^{\Im} (\mathbf{B})\) (where we mean the skew tensor product of graded algebras (Algebra, Chapter III).
\item
  Show that, if \(\mathbf{F}\) admits a complement in \(\mathbf{E}\) , the homomorphism defined in (a) is an isomorphism.
\item
   Take \(\mathbf{A} = \mathbf{Z}\) , \(\mathbf{E} = \mathbf{Z} / 4\mathbf{Z}\) , \(\mathbf{F} = 2\mathbf{Z} / 4\mathbf{Z}\) ; show that the homomorphism defined in (a) is not injective in this case.
\end{enumerate}

\begin{enumerate}
\def\labelenumi{\arabic{enumi}.}
\setcounter{enumi}{7}
\item
  Let A be a filtered ring, \((\mathbf{A}_{n})\) its filtration, E a filtered A-module and \((\mathbf{E}_{n})\) its filtration. Show that if \(A_{0} = A\) and \(E_{0} = E\) , the mapping \((a, x) \mapsto ax\) of \(A \times E\) to E is uniformly continuous with the topologies defined by the filtrations.
\item
  Give an example of two filtered rings A, B whose filtrations are exhaustive and separated and a non-surjective homomorphism \(u: A \to B\) compatible with the filtrations and such that \(\text{gr}(u)\) is bijective. From this deduce a counterexample to Corollary 1 to Proposition 12 when the ring A is not assumed to be complete and a counter-example to Proposition 13 when E is no longer assumed to be finitely generated (use Algebra, Chapter VIII, §7, Exercise 3(b)).
\item
  Let A be a Noetherian ring and \(\sigma\) an automorphism of A. Show that the ring E defined in Algebra, Chapter IV, § 5, Exercise 10 (b) is a (left or right) Noetherian ring.
\item
  Show that, if E is a vector space of dimension \(\geqslant2\) over a commutative field, the tensor algebra of E is a ring which is neither left nor right Noetherian (if a, b are two linearly independent vectors in E, consider the left ideal (or right ideal) generated by the elements \(a^{n}b^{n}\) for \(n\geqslant1\) ).
\end{enumerate}

 ¶ 12. Let K be a commutative field, A the polynomial ring \(K[X_{i}]_{i\in I}\) in an arbitrary infinite family of indeterminates and m the (maximal) ideal of A generated by the \(X_{i}\) . If we set \(A_{i}=A/m^{i+1}\) , the \(A_{i}\) and the canonical homomorphisms \(h_{ij}:A/m^{j+1}\to A/m^{i+1}\) for \(i\leqslant j\) satisfy the conditions of Proposition

14; the ring \(\lim A_{i}\) is the completion \(\hat{A}\) of \(A\) with respect to the m-adic topology and the kernel of the canonical homomorphism \(\hat{A} \to A_{i}\) is equal to \((m^{i})^{\wedge}\) , the closure of \(m^{i}\) in \(\hat{A}\) .

\begin{enumerate}
\def\labelenumi{(\alph{enumi})}
\item
  Show that \(\hat{\mathbf{A}}\) is canonically identified with the ring of formal power series in the \(\mathbf{X}_{\iota}\) , each of which has only a finite number of terms of given degree (Algebra, Chapter IV, § 5, Exercise 1).
\item
  From now on take \(\mathbf{I} = \mathbf{N}\) . Show that \(\hat{\mathfrak{m}} \neq \hat{\mathbf{A}}. \mathfrak{m}\) (consider the formal power series \(\sum_{n=1}^{\infty} \mathbf{X}_n^n\) ).
\item
  Suppose that K is a finite field. Show that \((\hat{\mathfrak{m}})^{2} \neq (\mathfrak{m}^{2})^{\bullet}\) . (Show first the following result: for every integer k \textgreater{} 0, there exists an integer \(n_{k}\) and for every integer \(n \geqslant n_{k}\) a homogeneous polynomial \(F_{n}\) of degree n in \(n^{2}\) indeterminates with coefficients in K such that \(F_{n}\) cannot be the sum of the terms of degree n in any polynomial of the form \(P_{1}Q_{1} + \cdots + P_{k}Q_{k}\) , where the \(P_{i}\) and \(Q_{i}\) are polynomials without constant term in the same \(n^{2}\) indeterminates).
\end{enumerate}

Deduce that \(\hat{A}\) is not complete with respect to the \(\hat{m}\) -adic topology.

\begin{enumerate}
\def\labelenumi{\arabic{enumi}.}
\setcounter{enumi}{12}
\tightlist
\item
  Let K be a field, \(A = K[[X]]\) the ring of formal power series and m its maximal ideal, so that A is Hausdorff and complete with the m-adic topology (no. 6, Corollary to Proposition 6). On the additive group A consider the filtration \((\mathrm{E}_{n})\) such that \(E_{0} = A\) and \(E_{n}\) is the intersection of \(m^{n}\) and the ring K{[}X{]}; this filtration is exhaustive and separated, the topology T which it defines on A is compatible with the additive group structure on A and is finer than the m-adic topology but A is not a complete group with the topology T (consider the sequence of polynomials \((1 - X^{n})/(1 - X)\) ).
\end{enumerate}

14 Let A be a ring (not necessarily commutative) and E a left A-module. A topology on E is called linear if it is invariant under translations and 0 admits a fundamental system of neighbourhoods which are submodules of E; E is then said to be linearly topologized. A linear topology on E is compatible with its additive group structure and defines on E a topological A-module structure if A is given the discrete topology. On any A-module the discrete topology and the coarsest topology are linear topologies.

\begin{enumerate}
\def\labelenumi{(\alph{enumi})}
\item
  If E is a linearly topologized module and F a submodule of E, the induce topology on F and the quotient topology of that on E by F are linear topologies. If \((\mathrm{E}_{\alpha}, f_{\alpha\beta})\) is an inverse system of linearly topologized A-modules, where the \(f_{\alpha\beta}\) are continuous linear mappings, the topological A-module \(E = \lim_{\leftarrow} E_{\alpha}\) is linearly topologized.
\item
  Let E be a linearly topologized A-module. There exists a fundamental system \((\mathrm{V}_{\lambda})_{\lambda \in L}\) of neighbourhoods of 0 in E consisting of open (and closed) submodules; if \(\phi_{\lambda\mu}: E/V_{\mu} \to E/V_{\lambda}\) is the canonical mapping when \(V_{\lambda} \supset V_{\mu}\) , the family \((\mathrm{E}/\mathrm{V}_{\lambda}, \phi_{\lambda\mu})\) is an inverse system of discrete A-modules; the topological A-module \(\tilde{E} = \lim_{\leftarrow} E/V_{\lambda}\) is identified with the Hausdorff completion of E, which is therefore linearly topologized (General Topology, Chapter III, § 7, no. 3).
\item
  A Hausdorff linearly topologized A-module E is discrete if it is Artinian or if there exists a least element in the set of submodules \(\neq0\) of E.
\end{enumerate}

\begin{enumerate}
\def\labelenumi{\arabic{enumi}.}
\setcounter{enumi}{14}
\tightlist
\item
  \begin{enumerate}
  \def\labelenumii{(\alph{enumii})}
  \tightlist
  \item
    Let E be a linearly topologized A-module (Exercise 14). For a filter base B on E consisting of linear varieties to admit a cluster point, it is necessary and sufficient that there exists a convergent filter base \(B' \supset B\) consisting of linear affine varieties. E is said to be linearly compact if it is Hausdorff and every filter base on E consisting of affine linear varieties admits at least one cluster point. Every Artinian module is linearly compact with the discrete topology. Every linearly compact submodule of a Hausdorff linearly topologized module F is closed in F.
  \end{enumerate}
\end{enumerate}

\begin{enumerate}
\def\labelenumi{(\alph{enumi})}
\setcounter{enumi}{1}
\item
  If E is a linearly compact A-module and u a continuous linear mapping of E to a Hausdorff linearly topologized A-module F, \(u(\mathbf{F})\) is a linearly compact submodule of F.
\item
  Let E be a Hausdorff linearly topologized A-module and F a closed submodule of E. For E to be linearly compact, it is necessary and sufficient that F and E/F be so.
\item
  Every product of linearly compact modules is linearly compact (consider a filter base which is maximal among the filter bases consisting of affine linear varieties.
\item
  Let \((\mathrm{E}_{\alpha}, f_{\alpha \beta})\) be an inverse system of linearly topologized modules relative to a directed indexing set; suppose that the \(f_{\alpha \beta}\) are continuous linear mappings and that, for \(\alpha \leqslant \beta\) , \(f_{\alpha \beta}^{-1}(0)\) is a linearly compact submodule of \(\mathbf{E}_{\beta}\) . Let \(\mathbf{E} = \varprojlim \mathbf{E}_{\alpha}\) and let \(f_{\alpha}\) be the canonical mapping \(\mathbf{E} \to \mathbf{E}_{\alpha}\) ; show that, for all \(\alpha, f_{\alpha}(\mathbf{E}) = \bigcap_{\alpha \leqslant \beta} f_{\alpha \beta}(\mathbf{E}_{\beta})\) (in particular, if the \(f_{\alpha \beta}\) are surjective, so are the \(f_{\alpha}\) ) and \(f_{\alpha}^{-1}(0)\) is linearly compact (use (d) and General Topology, Chapter I, Appendix, no. 2, Theorem 1).
\end{enumerate}

¶ 16. (a) Let E be a Hausdorff linearly topologized A-module. Show that the following properties are equivalent:

(α) E is linearly compact (Exercise 15).

(β) For every continuous linear mapping \(u\) from \(E\) to a Hausdorff linearly topologized A-module \(F\) , \(u(E)\) is a closed submodule of \(F\) .

\begin{enumerate}
\def\labelenumi{(\Alph{enumi})}
\setcounter{enumi}{24}
\tightlist
\item[(γ)]
  With every linear topology (Hausdorff or not) on E coarser than the given topology, E is complete.
\end{enumerate}

(δ) E is complete and there is a fundamental system \((\mathbf{U}_{\lambda})\) of open neighbourhoods of 0 in E, consisting of submodules and such that the \(E/U_{\lambda}\) are linearly compact discrete A-modules (cf.~§ 3, Exercise 5).

(To see that \((\gamma)\) implies \((\delta)\) consider an open submodule \(F\) of \(E\) and a filter base \(\mathfrak{B}\) on \(E / F\) consisting of affine linear varieties. For all \(V \in \mathfrak{B}\) , let \(M_V\) be the inverse image in E of the direction submodule of V in E/F; consider on E the linear topology in which the \(M_{V}\) form a fundamental system of neighbourhoods of 0.)

\begin{enumerate}
\def\labelenumi{(\alph{enumi})}
\setcounter{enumi}{1}
\item
  Let E be a Hausdorff linearly topologized A-module and M a linearly compact submodule of E. Show that, for every closed submodule F of E, \(M + F\) is closed in E (consider the image of M in E/F).
\item
  Let E be a linearly compact A-module and F a Hausdorff linearly topologized A-module. Show that for ever closed submodule M of \(E \times F\) , the projection of M onto F is closed. Obtain the converse (cf.~General Topology, Chapter I, § 10, no. 2).
\item
  Let E be a linearly compact A-module and u a continuous linear mapping of E to a Hausdorff linearly topologized A-module F. Show that, for every filter base B on E consisting of affine linear varieties, the image under u of the set of cluster points of B is the set of cluster points of \(u(\mathfrak{B})\) . In particular, for every closed submodule M of E, \(\bigcap_{N\in\mathfrak{B}}(M+N)=M+\bigcap_{N\in\mathfrak{B}}N\) .
\end{enumerate}

\begin{enumerate}
\def\labelenumi{\arabic{enumi}.}
\setcounter{enumi}{16}
\tightlist
\item
  A linear topology \(\mathcal{T}\) on an A-module E is called minimal if it is Hausdorff and there exists no Hausdorff linear topology strictly coarser than \(\mathcal{T}\) .
\end{enumerate}

\begin{enumerate}
\def\labelenumi{(\alph{enumi})}
\item
  For a Hausdorff linear topology T on E to be minimal, it is necessary and sufficient that every filter base B on E, consisting of affine linear varieties and having a single cluster point, be convergent to this point. (To see that the condition is necessary, observe that, when M runs through B and V a fundamental system of neighbourhoods of 0 consisting of submodules, the M + V form a filter base with the same cluster points as B. To see that the condition is sufficient, note that a filter base consisting of open submodules, whose intersection is reduced to 0, is a fundamental system of neighbourhoods of 0 in a Hausdorff linear topology coarser than T.)
\item
  For a discrete topology on an A-module E to be minimal, it is necessary and sufficient that E be Artinian.
\item
  If \(\mathcal{T}\) is minimal, the topology induced by \(\mathcal{T}\) on any closed submodule of \(\mathbf{E}\) is minimal.
\end{enumerate}

¶ 18. In an A-module E, a submodule M ≠ E is called sheltered if there exists a least element in the set of submodules ≠0 in E/M.

\begin{enumerate}
\def\labelenumi{(\alph{enumi})}
\item
  Show that every submodule \(N \neq E\) of E is an intersection of sheltered submodules (for all \(x \notin N\) , consider a maximal element in the set of submodules of E containing N and not containing x).
\item
  Let \(\mathcal{T}\) be a Hausdorff linear topology on E and let \(\mathcal{T}^*\) be the linear topology with fundamental system of neighbourhoods of 0 the filter base generated by the submodules of E which are open under \(\mathcal{T}\) and sheltered. Show that \(\mathcal{T}^*\) is Hausdorff and that every submodule which is closed under \(\mathcal{T}\) is also closed under \(\mathcal{T}^*\) (note that every submodule which is closed under \(\mathcal{T}\) is an intersection of submodules which are open under \(\mathcal{T}\) ). Deduce that, if E is complete with respect to \(\mathcal{T}^*\) , it is complete with respect to \(\mathcal{T}\) (General Topology, Chapter III, § 3, no. 5, Proposition 9).
\item
  Suppose that \(\mathcal{T}\) is linearly compact. Then show that \(\mathcal{T}^*\) is linearly compact and is the coarsest Hausdorff linear topology coarser than \(\mathcal{T}\) ; in particular \(\mathcal{T}^*\) is a minimal linear topology (Exercise 17). (Let \(\mathfrak{B}\) be a filter base consisting of submodules closed under \(\mathcal{T}\) whose intersection is 0; show that, for every submodule U which is open under \(\mathcal{T}\) and sheltered, there exists \(\mathbf{M} \in \mathfrak{B}\) such that \(\mathbf{M} \subset \mathbf{U}\) , using Exercise 16 (d)).
\item
  Let \(\mathbf{F}\) be a Hausdorff linearly topologized A-module and \(\mathcal{T}_1\) its topology; show that, if \(u: \mathrm{E} \to \mathrm{F}\) is a continuous linear mapping with respect to the topologies \(\mathcal{T}\) and \(\mathcal{T}_1\) , \(u\) is also continuous with respect to the topologies \(\mathcal{T}^*\) and \(\mathcal{T}_1^*\) .
\end{enumerate}

¶ 19. (a) Let E be a linearly compact A-module. Show that the following conditions are equivalent:

(α) For every continuous linear mapping u from E to a Hausdorff linearly topologized A-module F, u is a strict morphism (General Topology, Chapter III, § 2, no. 8) from E to F.

(β) For every closed submodule \(\mathbf{F}\) of \(\mathbf{E}\) , the quotient topology on \(\mathbf{E} / \mathbf{F}\) is a minimal topology (Exercise 17).

(γ) E is complete and there is a fundamental system \((\mathbf{U}_{\lambda})\) of open neighbourhoods of 0 in E, consisting of submodules and such that the \(E/U_{\lambda}\) are Artinian modules.

(To see that \((\gamma)\) implies \((\beta)\) , reduce it to the case where \(F = 0\) and use Exercises 17 and 18.)

If E satisfies the equivalent conditions \((\alpha)\) , \((\beta)\) and \((\gamma)\) , E is called strictly linearly compact.

\begin{enumerate}
\def\labelenumi{(\alph{enumi})}
\setcounter{enumi}{1}
\item
  Let E be a Hausdorff linearly topologized A-module and F a closed submodule of E. For E to be strictly linearly compact, it is necessary and sufficient that F and E/F be so.
\item
  Every inverse limit of strictly linearly compact modules is strictly linearly compact.
\item
  Let E be a Hausdorff linearly topologized A-module and u a linear mapping from E to a strictly linearly compact A-module F. Show that if the graph of u is closed in \(E \times F\) , u is continuous (use Exercise 16 (c) and the fact that, if M is closed in E and E/M is Artinian, M is open in E).
\end{enumerate}

¶ 20. (a) Show that a discrete linearly compact module cannot be the direct sum of an infinity of submodules which are not reduced to 0.

\begin{enumerate}
\def\labelenumi{(\alph{enumi})}
\setcounter{enumi}{1}
\item
  Give an example of a minimal linear topology (Exercise 17) which is not linearly compact (consider an infinite direct sum of simple modules no two of which are isomorphic).
\item
  Let E be a linearly compact module such that there exists a family of simple submodules whose sum is dense in E. Show that: (1) E is strictly linearly compact (apply (a) to the quotient of E by an open submodule); (2) E is isomorphic to the (topological) product of a family of discrete simple modules. (Consider the sets O of maximal open submodules of E such that for every finite
\end{enumerate}

\[
\mathrm{E} / \left(\bigcap_ {k = 1} ^ {n} \mathrm{G} _ {k}\right)
\]

sequence \((\mathbf{G}_k)_{1\leqslant k\leqslant n}\) of \(n\) distinct elements of \(\mathfrak{O}\) , \(\operatorname{E} / \left(\bigcap_{k=1}^{\infty}\mathbf{G}_k\right)\) is the direct sum of \(n\) simple submodules. Show that there exists a maximal set \(\mathfrak{O}\) (with respect to inclusion) and that the intersection of all the submodules \(\mathbf{G}\) belonging to such a set is reduced to 0; use the fact that \(\{0\}\) is an intersection of maximal open submodules and that for every maximal open submodules \(\mathbf{G}\) of \(\mathbf{E}\) there is at least one simple submodule of \(\mathbf{E}\) not contained in \(\mathbf{G}\) . Conclude by using the fact that \(\mathbf{E}\) is strictly linearly compact.)

\begin{enumerate}
\def\labelenumi{(\alph{enumi})}
\setcounter{enumi}{3}
\tightlist
\item
  Deduce from (c) that every linearly compact vector space over a field \(\mathbf{K}\) is strictly linearly compact and isomorphic to a product \(\mathbf{K}_s^{\mathrm{I}}\) .
\end{enumerate}

¶ 21. Recall that on a ring A a topology is called (left) linear if it is a linear topology on the A-module \(\mathbf{A}_s\) and it is compatible with the ring structure on A (Chapter II, § 2, Exercise 16). A topological ring A is called (left) linearly compact (resp. (left) strictly linearly compact) if \(\mathbf{A}_s\) is a linearly compact (resp. strictly linearly compact) A-module.

\begin{enumerate}
\def\labelenumi{(\alph{enumi})}
\item
  Show that, if A is left linearly compact with the topology T, the topology \(T^{*}\) defined in Exercise 18 (b) is also compatible with the ring structure on A (use Exercise 18 (d)).
\item
  Suppose that \(\mathbf{A}\) is commutative; let \(u \neq 0\) be an idempotent of \(\mathbf{A} / \Re\) , where \(\Re\) is the Jacobson radical of \(\mathbf{A}\) . Show that if \(\mathbf{A}\) is linearly compact with the discrete topology, there exists a unique idempotent \(e \in \mathbf{A}\) whose image in \(\mathbf{A} / \Re\) is equal to \(u\) . (Show that among the linear affine varieties \(x + \mathfrak{b}\) , where \(\mathfrak{b} \subset \Re\) and \(x^2 - x \in \mathfrak{b}\) , which are contained in the class \(u\) , there is a minimal element \(e + \mathfrak{a}\) . Deduce that \(e^2 = e\) and \(\mathfrak{a} = 0\) , by considering the element \(e^2 - e = r\) and showing that \(r \in \mathrm{Ar}^2\) , using Algebra, Chapter VIII, § 6, Exercise 10 (a). Prove the uniqueness of \(e\) by showing that if \(e_1, e_2\) are two idempotents of \(\mathbf{A}\) such that \(e_1 e_2 \in \Re\) , then \(e_1 e_2 = 0\) .)
\item
  Show that every commutative ring A which is linearly compact with the discrete topology is a direct composition of a finite number of local rings (which are linearly compact with the discrete topology). (Consider first the case, where \(\Re = 0\) , using Exercise 15 (b) and Chapter II, § 1, no. 2, Proposition 5, then apply (b) to the idempotents of A/ \(\Re\) .)
\item
  Show that every linearly compact (resp. strictly linearly compact) commutative ring A is the product of a family of linearly compact (resp. strictly linearly compact) local rings. (Let \((\mathfrak{m}_{\lambda})\) be the family of open maximal ideals of A; for every open ideal \(a\) not containing \(\mathfrak{m}_{\lambda}\) , let \(\mathrm{A}_{\lambda}^{\prime} \times \mathrm{A}_{\lambda}^{\prime \prime}\) be the decomposition of A/a as a direct composition of two rings such that \(\mathrm{A}_{\lambda}^{\prime}\) is the local ring with maximal ideal \((\mathfrak{m}_{\lambda} + \mathfrak{a}) / a\) defined in (c); let \(e_{\lambda}^{\prime}(a)\) be the unit element of \(\mathrm{A}_{\lambda}^{\prime}\) considered as a class mod. \(a\) in A; the \(e_{\lambda}^{\prime}(a)\) form a filter base which converges in A to an idempotent \(e_{\lambda}\) and \(Ae_{\lambda}\) is a closed ideal of A of which \(e_{\lambda}\) is the unit element. Show that the ring \(Ae_{\lambda}\) is a local ring and that A is isomorphic to the product of the \(Ae_{\lambda}\) .
\end{enumerate}

¶ 22. (a) Let A be a local ring and m its maximal ideal. Show that, if, for a linear topology T on A, A is strictly linearly compact, T is coarser than the m-adic topology.

\begin{enumerate}
\def\labelenumi{(\alph{enumi})}
\setcounter{enumi}{1}
\item
  Let A be a commutative ring and m an ideal of A. For A to be strictly linearly compact with the m-adic topology, it is necessary and sufficient that it be Hausdorff and complete with this topology, that A/m be an Artinian ring and \(m/m^{2}\) and (A/m)-module of finite length * (in other words, A is a complete Noetherian semi-local ring) *.
\item
  Let I be an infinite indexing set, K a finite field and \(A = K[[X_{t}]]_{t \in I}\) the algebra of formal power series in the family of indeterminates \((\mathbf{X}_{t})_{t \in I}\) (Algebra, Chapter IV, § 5, Exercise 1); A is a local ring. As a vector space over K, A may be identified with the product space \(\mathbf{K}^{\mathbf{N}^{(1)}}\) ; if K is given the discrete topology and A the product topology T, show that T is a linear topology on the ring A with which A is strictly linearly compact; if m is the maximal ideal of A, show, by arguing as in Exercise 12 (c), that A is not complete with the m-adic topology.
\end{enumerate}

¶ 23. Let A be a commutative ring with a linear topology T.

\begin{enumerate}
\def\labelenumi{(\alph{enumi})}
\item
  Show that the ideals \(\mathfrak{a}\) of \(\mathbf{A}\) , which are closed under \(\mathcal{T}\) and such that \(\mathbf{A} / \mathfrak{a}\) is Artinian, form a fundamental system of neighbourhoods of 0 under a linear topology \(\mathcal{T}^{(c)}\) on \(\mathbf{A}\) , which is coarser than \(\mathcal{T}\) (use Algebra, Chapter VIII, §2, Exercise 12); then \((\mathcal{T}^{(c)})^{(c)} = \mathcal{T}^{(c)}\) . Let \(\mathbf{B}\) be the Hausdorff completion ring of \(\mathbf{A}\) with respect to the topology \(\mathcal{T}^{(c)}\) ; show that \(\mathbf{B}\) is strictly linearly compact; \(\mathbf{B}\) is said to be the strictly linearly compact ring associated with \(\mathbf{A}\) .
\item
  Let \(\mathcal{T}_{\omega}(\mathbf{A})\) be the discrete topology on \(\mathbf{A}\) , \(\mathcal{T}_c(\mathbf{A})\) the topology \((\mathcal{T}_{\omega}(\mathbf{A}))^{(c)}\) . For the topology \(\mathcal{T}^{(c)}\) on \(\mathbf{A}\) to be Hausdorff, it is necessary and sufficient that \(\mathcal{T}\) be Hausdorff and that, for every ideal \(\mathfrak{b}\) of \(\mathbf{A}\) open under \(\mathcal{T}\) , the topology \(\mathcal{T}_c(\mathbf{A}/\mathfrak{b})\) be Hausdorff. * If \(\mathbf{A}\) is the ring of a non-discrete valuation of height 1 (Chapter VI), \(\mathcal{T}_c(\mathbf{A})\) is not Hausdorff. *
\item
  Suppose that A is linearly compact under \(\mathcal{T}\) ; then A is complete (but in general not Hausdorff) with \(\mathcal{T}^{(c)}\) . Then for \(\mathcal{T}^{(c)}\) to be Hausdorff, it is necessary and sufficient that there exist on A a linear topology which is strictly linearly compact and coarser than \(\mathcal{T}\) ; \(\mathcal{T}^{(c)}\) is then the unique topology with these properties.
\item
  Let E be a linearly topologized and strictly linearly compact A-module. Show that, if A is given the topology \(\mathcal{T}_{c}(A)\) , E is a topological ring (note that if F is an Artinian A-module, then, for all \(x \in F\) , the annihilator a of x is such that A/a is Artinian). Deduce that, if B is the Hausdorff completion of A with respect to the topology \(\mathcal{T}_{c}(A)\) , E can be considered as a topological B-module; the ring B being isomorphic to a product of strictly linearly compact local rings
\end{enumerate}

\(B_{\lambda}\) (Exercise 21 (d)), show that E is isomorphic to a product of strictly linearly compact submodules \(E_{\lambda}\) , where \(E_{\lambda}\) is annihilated by the \(B_{\mu}\) of index \(\mu \neq \lambda\) and can therefore be considered as a topological \(B_{\lambda}\) -module (if \(e_{\lambda}\) is the unit element of B, take \(E = e_{\lambda}.E\) ).

\begin{enumerate}
\def\labelenumi{\arabic{enumi}.}
\setcounter{enumi}{23}
\tightlist
\item
  On a commutative ring A let \(\mathcal{T}_{m}(A)\) denote the linear topology which has a fundamental system of neighbourhoods of 0 the products (equal to the intersections) of a finite number of powers of maximal ideals (cf.~no. 13, Proposition 17); if every finite intersection of ideals \(\neq0\) of A is \(\neq0\) , let \(\mathcal{T}_{u}(A)\) denote the linear topology for which a fundamental system of neighbourhoods of 0 consists of all the ideals \(\neq\{0\}\) of A.
\end{enumerate}

\begin{enumerate}
\def\labelenumi{(\alph{enumi})}
\item
  Show that \(\mathcal{T}_c(\mathbf{A})\) (Exercise 23) is coarser than \(\mathcal{T}_m(\mathbf{A})\) (observe that the Jacobson radical of an Artinian ring is nilpotent); give an example where \(\mathcal{T}_c(\mathbf{A}) \neq \mathcal{T}_m(\mathbf{A})\) (cf.~Exercise 22 (c)). * If A is the ring of a non-discrete valuation of height 1 (Chapter VI), then \(\mathcal{T}_c(\mathbf{A}) = \mathcal{T}_m(\mathbf{A})\) and \(\mathcal{T}_u(\mathbf{A}) \neq \mathcal{T}_m(\mathbf{A})\) . *
\item
  The topology \(\mathcal{T}_m(A)\) is the coarsest linear topology on \(A\) under which the maximal ideals of \(A\) are closed and every power of an open ideal is an open ideal (note that under a linear topology on \(A\) , under which a maximal ideal \(m\) is closed, \(m\) is necessarily open).
\item
  If A is a Noetherian ring, then \(\mathcal{T}_{m}(\mathbf{A}) = \mathcal{T}_{c}(\mathbf{A})\) (observe that, if \(m\) is a maximal ideal of A, \(\mathbf{A} / m^{k}\) is an Artinian ring for every integer \(k \geqslant 1\) , noting that \(\mathfrak{m}^h / \mathfrak{m}^{h+1}\) is an (A/m)-module which is necessarily of finite length). Give an example of a Noetherian local ring A for which \(\mathcal{T}_{m}(\mathbf{A}) \neq \mathcal{T}_{u}(\mathbf{A})\) (consider a local ring of a polynomial ring over a field).
\end{enumerate}

\begin{enumerate}
\def\labelenumi{\arabic{enumi}.}
\setcounter{enumi}{24}
\item
  Let A be a topological ring and E a finitely generated left A-module. Show that there exists on E a topology (called canonical) compatible with its A-module structure and finer than all the others (write E in the form \(A_{s}^{n}/R\) , give \(A_{s}^{n}\) the product topology and E the quotient topology). If \(E'\) is a topological A-module and \(u: E \to E'\) an A-linear mapping, show that u is continuous with the canonical topology on E. If further \(E'\) is finitely generated and has the canonical topology, u is a strict morphism. If the topology on A is defined by a filtration ( \(A_{n}\) ), show that the canonical topology on E is defined by the filtration ( \(A_{n}E\) ).
\item
  Let A be a commutative ring, m an ideal of A, \((a_{\lambda})_{\lambda \in L}\) a system of generators of m and \(\hat{A}\) the Hausdorff completion of A with respect to the m-adic topology.
\end{enumerate}

Show that, if \(\mathbf{A}'\) denotes the ring of formal power series in a family \((\mathrm{T}_{\lambda})_{\lambda \in \mathbf{L}}\) of indeterminates with only a finite number of terms of given degree (Exercise 12), \(\hat{\mathbf{A}}\) is isomorphic to the quotient of \(\mathbf{A}'\) by the closure \(\mathfrak{b}\) of the ideal of \(\mathbf{A}'\) generated by the \(T_{\lambda} - a_{\lambda}\) (use Theorem 1 of no. 8). In particular, the ring \(Z_{p}\) of p-adic integers is isomorphic to the quotient \(\mathbf{Z}[[T]]/(T-p)\) .

¶ 27. (a) Let A be a commutative ring with a linear topology. For every multiplicative subset S of A, let \(\mathrm{A}\{\mathrm{S}^{-1}\}\) denote the Hausdorff completion of the ring \(\mathrm{S}^{-1}\mathrm{A}\) with the topology for which a fundamental system of neighbourhoods of 0 consists of the ideals \(\mathrm{S}^{-1}\mathrm{U}_{\lambda}\) , where \((\mathrm{U}_{\lambda})\) is a fundamental system of neighbourhoods of 0 in A consisting of ideals of A. Show that \(\mathrm{A}\{\mathrm{S}^{-1}\}\) is canonically isomorphic to the inverse limit of the rings \(\mathrm{S}_{\lambda}^{-1}(\mathrm{A}/\mathrm{U}_{\lambda})\) where \(\mathrm{S}_{\lambda}\) is the canonical image of S in \(\mathrm{A}/\mathrm{U}_{\lambda}\) . If \(\hat{\mathrm{A}}\) is the Hausdorff completion of A, \(\mathrm{A}\{\mathrm{S}^{-1}\}\) is canonically isomorphic to \(\hat{\mathrm{A}}\{\mathrm{S}'^{-1}\}\) , where \(\mathrm{S}'\) is the canonical image of S in \(\hat{\mathrm{A}}\) .

\begin{enumerate}
\def\labelenumi{(\alph{enumi})}
\setcounter{enumi}{1}
\item
  For \(A\{S^{-1}\}\) to be reduced to 0, it is necessary and sufficient that, in A, 0 belong to the closure of S. Deduce an example where A is Hausdorff and complete but \(S^{-1}A\) is not so with the topology defined in (a).
\item
  For every continuous homomorphism \(u\) of \(A\) to a complete Hausdorff linearly topologized ring \(B\) such that \(u(S)\) consists of invertible elements of \(B\) , show that \(u = u' \circ j\) , where \(j: A \to A\{S^{-1}\}\) is the canonical mapping and \(u'\) is continuous; moreover \(u'\) is determined uniquely.
\item
  Let \(S_{1}\) , \(S_{2}\) be two multiplicative subsets of A and let \(S_{2}^{\prime}\) be the canonical image of \(S_{2}\) in \(A\{S^{-1}\}\) ; define a canonical isomorphism of \(A\{(S_{1}S_{2})^{-1}\}\) onto \(A\{S_{1}^{-1}\}\{S_{2}'^{-1}\}\) .
\item
  Let \(\mathfrak{a}\) be an open ideal of A and let \(\mathfrak{a}\{\mathrm{S}^{-1}\}\) be the Hausdorff completion of \(\mathrm{S}^{-1}\mathfrak{a}\) with respect to the topology induced by that on \(\mathrm{S}^{-1}\mathrm{A}\) ; show that \(\mathfrak{a}\{\mathrm{S}^{-1}\}\) is canonically identified with an open ideal of \(\mathrm{A}\{\mathrm{S}^{-1}\}\) and that the discrete ring \(\mathrm{A}\{\mathrm{S}^{-1}\} / \mathfrak{a}\{\mathrm{S}^{-1}\}\) is isomorphic to \(\mathrm{S}^{-1}(\mathrm{A} / \mathfrak{a})\) . Conversely, if \(\mathfrak{a}'\) is an open ideal of \(\mathrm{A}\{\mathrm{S}^{-1}\}\) , its inverse image \(\mathfrak{a}\) in A is an open ideal such that \(\mathfrak{a}' = \mathfrak{a}\{\mathrm{S}^{-1}\}\) . In particular, the mapping \(\mathfrak{p} \mapsto \mathfrak{p}\{\mathrm{S}^{-1}\}\) is an increasing bijection of the set of open prime ideals of A not meeting S onto the set of open prime ideals of \(\mathrm{A}\{\mathrm{S}^{-1}\}\) .
\item
  Let \(p\) be an open prime ideal of \(A\) and let \(S = A - p\) . Show that \(A\{S^{-1}\}\) is a local ring whose residue field is isomorphic to the field of fractions of \(A / p\) .
\item
  For every element \(f \in \mathbf{A}\) , let \(\mathbf{A}_{(f)}\) denote the ring \(\mathbf{A}\{\mathbf{S}_f^{-1}\}\) , where \(\mathbf{S}_f\) is the multiplicative set of the \(f^n\) ( \(n \geqslant 0\) ). If \(f\) runs through a multiplicative subset \(\mathbf{S}\) of \(\mathbf{A}\) , the \(\mathbf{A}_{(f)}\) form a direct system of rings whose direct limit is denoted by \(\mathbf{A}_{(\mathbf{S})}\) (without any topology); define a canonical homomorphism
\end{enumerate}

\[
\mathrm{A} _ {\{\mathrm{S} \}} \rightarrow \mathrm{A} \{\mathrm{S} ^ {- 1} \}.
\]

If \(S = A - p\) , where \(p\) is an open prime ideal of \(A\) , show that \(A_{(S)}\) is a local ring, the canonical homomorphism \(A_{(S)} \to A\{S^{-1}\}\) is local and the residue fields of \(A_{(S)}\) and \(A\{S^{-1}\}\) are canonically isomorphic (cf.~Chapter II, § 3, Exercise 16).

\begin{enumerate}
\def\labelenumi{(\alph{enumi})}
\setcounter{enumi}{7}
\tightlist
\item
  Suppose that the topology on A is the m-adic topology for some ideal m of A, that A is Hausdorff and complete and that \(m/m^{2}\) is a finitely generated (A/m)-module. Show that, if \(m' = m\{S^{-1}\}\) , the topology on \(A' = A\{S^{-1}\}\) is the \(m'\) -adic topology, \(m' = mA'\) and \(m'/m'^{2}\) is a finitely generated ( \(A'/m'\) )-module (use Proposition 14 of no. 10). If A is Noetherian, so is \(A'\) .
\end{enumerate}

\begin{enumerate}
\def\labelenumi{\arabic{enumi}.}
\setcounter{enumi}{27}
\tightlist
\item
  Let A be a commutative ring with a linear topology and E, F two topological A-modules which are linearly topologized. If V (resp. W) runs through the set of open submodules of E (resp. F), the submodules
\end{enumerate}

\[
\operatorname{Im} (\mathrm{V} \otimes_ {\mathbf {A}} \mathrm{F}) + \operatorname{Im} (\mathrm{E} \otimes_ {\mathbf {A}} \mathrm{W})
\]

of \(\mathbf{E} \otimes_{\mathbf{A}} \mathbf{F}\) form a fundamental system of neighbourhoods of 0 in \(\mathbf{E} \otimes_{\mathbf{A}} \mathbf{F}\) for a topology which is compatible with its module structure over the topological ring A and is called the tensor product of the given topologies on E and F. The Hausdorff completion \((\mathbf{E} \otimes_{\mathbf{A}} \mathbf{F})^{\wedge}\) of this A-module is an \(\hat{\mathbf{A}}\) -module called the completed tensor product of E and F.

\begin{enumerate}
\def\labelenumi{(\alph{enumi})}
\tightlist
\item
  Show that, if \((\mathrm{V}_{\lambda})\) (resp. \((\mathrm{W}_{\mu})\) ) is a fundamental system of neighbourhoods of 0 in E (resp. F) consisting of submodules, \((\mathbf{E} \otimes_{\mathbf{A}} \mathbf{F})^{\wedge}\) is canonically isomorphic to the inverse limit of the inverse system of A-modules
\end{enumerate}

\[
\left(\mathrm{E} / \mathrm{V} _ {\lambda}\right) \otimes_ {\mathrm{A}} \left(\mathrm{F} / \mathrm{W} _ {\mu}\right);
\]

deduce that \((\mathbf{E} \otimes_{\mathbf{A}} \mathbf{F})^{\wedge}\) is an \(\hat{\mathbf{A}}\) -module canonically isomorphic to \((\hat{\mathbf{E}} \otimes_{\hat{\mathbf{A}}} \hat{\mathbf{F}})\) ; it is also denoted by \(\hat{\mathbf{E}} \widehat{\otimes}_{\hat{\mathbf{A}}} \hat{\mathbf{F}}\) .

\begin{enumerate}
\def\labelenumi{(\alph{enumi})}
\setcounter{enumi}{1}
\item
  Let \(\mathbf{E}'\) , \(\mathbf{F}'\) be two topological A-modules which are linearly topologized and \(u: \mathbf{E} \to \mathbf{E}', v: \mathbf{F} \to \mathbf{F}'\) two continuous A-linear mappings; show that \(u \otimes v: \mathbf{E} \otimes \mathbf{F} \to \mathbf{E}' \otimes \mathbf{F}'\) is continuous with the tensor product topologies on \(\mathbf{E} \otimes \mathbf{F}\) and \(\mathbf{E}' \otimes \mathbf{F}'\) ; we denote by \(u \hat{\otimes} v\) the continuous linear mapping \((\mathbf{E} \otimes \mathbf{F})^{\wedge} \to (\mathbf{E}' \otimes \mathbf{F}')^{\wedge}\) corresponding to \(u \otimes v\) .
\item
  Let B, C be two commutative A-algebras with linear topologies such that the canonical mappings \(\mathbf{A} \to \mathbf{B}\) , \(\mathbf{A} \to \mathbf{C}\) are continuous (to abbreviate, B, C are called topological A-algebras). Show that \((\mathbf{B} \otimes_{\mathbf{A}} \mathbf{C})^{\wedge}\) has a canonical topological A-algebra structure called the completed tensor product of the algebras B and C. Define canonical continuous representations
\end{enumerate}

\[
\rho \colon \mathbf {B} \rightarrow (\mathbf {B} \otimes_ {\mathbf {A}} \mathbf {C}) ^ {\wedge}, \quad \sigma \colon \mathbf {C} \rightarrow (\mathbf {B} \otimes_ {\mathbf {A}} \mathbf {C}) ^ {\wedge}
\]

with the following property: for every complete Hausdorff commutative topological A-algebra D and every ordered pair of continuous A-homomorphisms \(u\colon B\to D\) , \(v\colon C\to D\) , there exists a unique continuous A-homomorphism \(w\colon(B\otimes_{A}C)^{\wedge}\to D\) such that \(u=w\circ\rho\) and \(v=w\circ\sigma\) .

\begin{enumerate}
\def\labelenumi{(\alph{enumi})}
\setcounter{enumi}{3}
\tightlist
\item
  Let m be an ideal of A such that the topology on A is the m-adic topology. If E and F are each given the m-adic topology, show that on \(E \otimes_{A} F\) the tensor product of the topologies on E and F is the m-adic topology. Deduce that \((\mathbf{E} \otimes_{\mathbf{A}} \mathbf{F})^{\wedge}\) is isomorphic to the inverse limits \(\varprojlim ((\mathrm{E} / \mathfrak{m}^{n+1}) \otimes_{\mathbf{A}} \mathbf{F})\) and \(\varprojlim (\mathbf{E} \otimes_{\mathbf{A}} (\mathbf{F} / \mathfrak{m}^{n+1}\mathbf{F}))\) .
\end{enumerate}

\begin{enumerate}
\def\labelenumi{\arabic{enumi}.}
\setcounter{enumi}{28}
\item
  Let A be a commutative ring, m an ideal of A such that \(\bigcap_{n\geq0}m^{n}=0\) and c is an element of A which is not a divisor of 0. Show that the transporters \(q_{n}=m^{n}:Ac\) are open under the m-adic topology and have zero intersection. If A is strictly linearly compact with the m-adic topology (Exercise 22(b)), \((q_{n})\) is a fundamental system of neighbourhoods of 0 in A for this topology.
\item
  Let K be an infinite commutative field and A the ring of formal power series \(K[[X,Y]]\) in two indeterminates, which is a complete Hausdorff Noetherian local ring. Define a multiplicative subset S of A such that \(S^{-1}A\) is not a semi-local ring (if \((\lambda_{n})\) is an infinite sequence of distinct elements of K, consider the prime ideals \(p_{n}=A(X+\lambda_{n}Y)\) of A).
\item
  Show that if A is a semi-local ring, so is the ring of formal power series A{[}{[}X{]}{]} (consider the Jacobson radical of this ring).
\end{enumerate}

